% GENERATED FILE. Edit canonical lessons, then run npm run generate:books.
% canonical-source-hash: fnv1a64:387b8484c1732161
% canonical-lessons: FR-C383-exigence, FR-C383-critere, FR-C383-eligible, FR-C383-admission, FR-C383-exclure

\chapter{Qui peut participer}
\label{ch:fr-qui-peut-participer}
\hlchaptermodality{383}

\begin{chapteropening}
\textbf{By the end of this chapter:} \emph{I can say who can take part: the requirements and criteria, being eligible, admission, and being left out.}

Complete the last lesson of chapter 383: I can say who can take part: the requirements and criteria, being eligible, admission, and being left out.
\end{chapteropening}

\section[l'exigence]{l'exigence — the requirement, the demand}

\label{lesson:FR-C383-exigence}



\begin{quote}
Before the new one: say the French for the penalty, then the French for to penalise.
\end{quote}

\subsection*{You'll want to know: l'exigence}
\textbf{l'exigence} — \textquotedblleft{}the requirement, the demand\textquotedblright{}.

\begin{quote}
Ce poste a beaucoup d'exigences.
\end{quote}

\begin{quote}
\emph{This job has a lot of requirements.}
\end{quote}

Feminine. \emph{Être exigeant} — to be demanding.

\begin{cousinweb}
\emph{exigence} is usually traced to Late Latin \emph{exigentia} \textquotedblleft{}urgency, need\textquotedblright{}, from \emph{exigere} \textquotedblleft{}to drive out, to demand\textquotedblright{} (Lewis \& Short), built on \emph{agere} \textquotedblleft{}to drive\textquotedblright{}.
\end{cousinweb}

\begin{grammarlens}[title={What you've built}]
One more word for telling, explaining and describing at B1.
\end{grammarlens}

\subsection*{Your turn}
\begin{itemize}
\raggedright
  \item \emph{Say it:} \emph{l'exigence}
  \item \emph{Say it:} \emph{l'exigence}, once more
  \item \emph{Say it:} a sentence of your own with \emph{l'exigence}
  \item \emph{Recall:} say the French for the penalty, then the French for to penalise, then say \emph{l'exigence} again
\end{itemize}

\begin{tcolorbox}[breakable,colback=teal!4,colframe=teal!35!black,title={Before you move on}]
What does l'exigence mean? (\textquotedblleft{}The requirement, the demand\textquotedblright{}.) Say it once more.
\end{tcolorbox}
\section[le critère]{le critère — the criterion}

\label{lesson:FR-C383-critere}



\begin{quote}
Before the new one: say the French for to penalise, then the French for the requirement.
\end{quote}

\subsection*{You'll want to know: le critère}
\textbf{le critère} — \textquotedblleft{}the criterion\textquotedblright{}.

\begin{quote}
Le prix est mon premier critère.
\end{quote}

\begin{quote}
\emph{Price is my first criterion.}
\end{quote}

The plural is \emph{les critères}.

\begin{cousinweb}
\emph{critère} is usually traced, as a learned word, to Late Latin \emph{critērium}, from Greek \emph{krit\'{\={e}}rion} \textquotedblleft{}a means of judging\textquotedblright{}, built on \emph{krínein} \textquotedblleft{}to judge\textquotedblright{}.
\end{cousinweb}

\begin{grammarlens}[title={What you've built}]
One more word for telling, explaining and describing at B1.
\end{grammarlens}

\subsection*{Your turn}
\begin{itemize}
\raggedright
  \item \emph{Say it:} \emph{le critère}
  \item \emph{Say it:} \emph{le critère}, once more
  \item \emph{Say it:} a sentence of your own with \emph{le critère}
  \item \emph{Recall:} say the French for to penalise, then the French for the requirement, then say \emph{le critère} again
\end{itemize}

\begin{tcolorbox}[breakable,colback=teal!4,colframe=teal!35!black,title={Before you move on}]
What does le critère mean? (\textquotedblleft{}The criterion\textquotedblright{}.) Say it once more.
\end{tcolorbox}
\section[éligible]{éligible — eligible}

\label{lesson:FR-C383-eligible}



\begin{quote}
Before the new one: say the French for the requirement, then the French for the criterion.
\end{quote}

\subsection*{You'll want to know: éligible}
\textbf{éligible} — \textquotedblleft{}eligible\textquotedblright{}.

\begin{quote}
À dix-huit ans, on est éligible au conseil municipal.
\end{quote}

\begin{quote}
\emph{At eighteen, you are eligible for the town council.}
\end{quote}

One form for both genders. \emph{Éligible à une aide} — eligible for aid.

\begin{cousinweb}
\emph{éligible} is usually traced to Late Latin \emph{ēligibilis} \textquotedblleft{}that may be chosen\textquotedblright{}, from \emph{ēligere} \textquotedblleft{}to pick out\textquotedblright{}, built on \emph{legere} \textquotedblleft{}to choose, to gather\textquotedblright{} (Lewis \& Short).
\end{cousinweb}

\begin{grammarlens}[title={What you've built}]
One more word for telling, explaining and describing at B1.
\end{grammarlens}

\subsection*{Your turn}
\begin{itemize}
\raggedright
  \item \emph{Say it:} \emph{éligible}
  \item \emph{Say it:} \emph{éligible}, once more
  \item \emph{Say it:} a sentence of your own with \emph{éligible}
  \item \emph{Recall:} say the French for the requirement, then the French for the criterion, then say \emph{éligible} again
\end{itemize}

\begin{tcolorbox}[breakable,colback=teal!4,colframe=teal!35!black,title={Before you move on}]
What does éligible mean? (\textquotedblleft{}Eligible\textquotedblright{}.) Say it once more.
\end{tcolorbox}
\section[l'admission]{l'admission — the admission, the entry}

\label{lesson:FR-C383-admission}



\begin{quote}
Before the new one: say the French for the criterion, then the French for eligible.
\end{quote}

\subsection*{You'll want to know: l'admission}
\textbf{l'admission} — \textquotedblleft{}the admission, the entry\textquotedblright{}.

\begin{quote}
Les résultats d'admission arrivent en juin.
\end{quote}

\begin{quote}
\emph{The admission results come out in June.}
\end{quote}

Feminine. From \emph{admettre}.

\begin{cousinweb}
\emph{admission} is a learned borrowing from Latin \emph{admissiō} \textquotedblleft{}a letting in\textquotedblright{} (EtymDB), from \emph{admittere} \textquotedblleft{}to let in\textquotedblright{}, built on \emph{mittere} \textquotedblleft{}to send\textquotedblright{} (Lewis \& Short).
\end{cousinweb}

\begin{grammarlens}[title={What you've built}]
One more word for telling, explaining and describing at B1.
\end{grammarlens}

\subsection*{Your turn}
\begin{itemize}
\raggedright
  \item \emph{Say it:} \emph{l'admission}
  \item \emph{Say it:} \emph{l'admission}, once more
  \item \emph{Say it:} a sentence of your own with \emph{l'admission}
  \item \emph{Recall:} say the French for the criterion, then the French for eligible, then say \emph{l'admission} again
\end{itemize}

\begin{tcolorbox}[breakable,colback=teal!4,colframe=teal!35!black,title={Before you move on}]
What does l'admission mean? (\textquotedblleft{}The admission, the entry\textquotedblright{}.) Say it once more.
\end{tcolorbox}
\section[exclure]{exclure — to exclude, to rule out}

\label{lesson:FR-C383-exclure}



\begin{quote}
Before the new one: say the French for eligible, then the French for the admission.
\end{quote}

\subsection*{You'll want to know: exclure}
\textbf{exclure} — \textquotedblleft{}to exclude, to rule out\textquotedblright{}.

\begin{quote}
On ne peut pas exclure un retard.
\end{quote}

\begin{quote}
\emph{A delay cannot be ruled out.}
\end{quote}

Like \emph{conclure}: past participle \emph{exclu}.

\begin{cousinweb}
\emph{exclure} is a learned borrowing from Latin \emph{exclūdere} \textquotedblleft{}to shut out\textquotedblright{} (EtymDB; Lewis \& Short), built from \emph{ex} \textquotedblleft{}out\textquotedblright{} and \emph{claudere} \textquotedblleft{}to shut\textquotedblright{}.
\end{cousinweb}

\begin{grammarlens}[title={What you've built}]
One more word for telling, explaining and describing at B1.
\end{grammarlens}

\subsection*{Your turn}
\begin{itemize}
\raggedright
  \item \emph{Say it:} \emph{exclure}
  \item \emph{Say it:} \emph{exclure}, once more
  \item \emph{Say it:} a sentence of your own with \emph{exclure}
  \item \emph{Recall:} say the French for eligible, then the French for the admission, then say \emph{exclure} again
\end{itemize}

\begin{tcolorbox}[breakable,colback=teal!4,colframe=teal!35!black,title={Before you move on}]
What does exclure mean? (\textquotedblleft{}To exclude, to rule out\textquotedblright{}.) Say it once more.
\end{tcolorbox}
